% from main.tex
%%%%%%%%%%  scan idx 125  |  folio 112  %%%%%%%%%%
\textbf{Exercice 2.16.2.} Soient \(S\) un topos, et \(A\to A'\) un morphisme d'Anneaux de \(S\). On fait l'une des hypothèses suivantes :
\begin{enumerate}[label=(\roman*)]
\item \(A\to A'\) est surjectif de noyau un Idéal nilpotent
\item le foncteur \(A'\otimes_A-\) sur la catégorie des \(A\)-modules à gauche est exact et fidèle.
\end{enumerate}
Montrer que \(E\in\Ob\Dpx^-(_A S)\) est \(n\)-pseudo-cohérent (resp. pseudo-cohérent) si et seulement si \(A'\otimes_A^{\mathrm L} E\in\Ob\Dpx^-(_{A'}S)\) est \(n\)-pseudo-cohérent (resp. pseudo-cohérent).

Indications : Montrer d'abord qu'un \(A\)-Module à gauche \(L\) est de type fini si et seulement si \(A'\otimes_A L\) l'est (cf. Bourbaki, Alg. commutative, chap. I, § 3, prop. 11 et chap. II, § 3, prop. 4), puis se ramener à ce cas par troncature, en utilisant (2.5) et (2.10 b)).

\section*{3. Lien avec la notion classique de cohérence}

\textbf{3.0.} Les hypothèses et les notations sont celles de (2.0), à cela près qu'on ne suppose plus \(\mathcal C_0\) localement quasi-stable par noyau d'épimorphisme.

\textbf{Définition 3.1.} \emph{Soit \(F\in\Ob\mathcal C_S\). Nous dirons que \(F\) est précohérent (relativement à \(\mathcal C_0\))\footnote{Nous introduisons cette notion ici à titre d'intermédiaire technique, la seule notion à retenir est celle de cohérence.} si, pour tout \(U\) au-dessus de \(S\) et tout morphisme \(u:E\to F|U\) où \(E\in\Ob\mathcal C_{0,U}\), \(\Ker(u)\) est de \(\mathcal C_0\)-type fini (1.2). Nous dirons que \(F\) est cohérent (relativement à \(\mathcal C_0\)) si \(F\) est précohérent et de type fini.}

\textbf{Remarque 3.2.} Soit \(F\in\Ob\mathcal C_S\). Il est clair, sur les définitions, que si \(F\) est précohérent (resp. de type fini) alors tout sous-objet (resp. quotient) de \(F\) est précohérent (resp. de type fini).

\textbf{Lemme 3.3.} \textbf{a)} Soit
\[
  0\longrightarrow F\longrightarrow G\longrightarrow H
\]
